% One-time imported draft. Editable native source; no Markdown regeneration.
\section{Labels and the instruction
sequence}\label{001-console-calculation-p03-labels-and-the-instruction-sequence}

The assembler maintains a mapping from labels to locations. A data label
can name the first byte of a message; a text label can name an
instruction location. The colon introduces the name in both cases.
Executing a program does not execute the label as an additional
operation.

\texttt{.globl\ main} makes the symbol \texttt{main} available as an
exported entry name. The starter places its body at \texttt{main},
followed by the supplied exit sequence. With no branches in the body,
execution reaches these instructions in order.

The source layout separates declarations from executable operations. The
\texttt{.data} and \texttt{.text} directives tell the assembler where
subsequent material belongs. These directives are instructions to the
assembler; they are not processor instructions that print text or
calculate values at run time.
